\chapter*{Pengantar}
\addcontentsline{toc}{chapter}{Pengantar}

Alhamdulillaahiroobil 'aalamiin. Segala puji bagi Tuhan, Rabb semesta alam, yang telah menambahkan ilmu dan merezekikan kefahaman sehingga kami bisa menyelesaikan buku ini: Konsep Getaran dan implementasinya. Buku ini kami susun sebagai dokumentasi atas apa yang telah kami ajarkan di kelas agar dapat dipelajari lebih lanjut oleh para siswa dan sebagai usaha mengamalkan ilmu yang telah diberikan kepada kami.

Buku ini dirangkum secara gado-gado dari belajar dan mengajar beberapa mata kuliah: matematika, matematika rekayasa, sistem linear, dan pemrosesan sinyal. Secara umum, kesemua materi adalah dasar Matematika \textit{untuk} Rekayasa. Jadi, isinya mungkin tidak seperti yang anda harapkan. Tapi, meski begitu, kami berharap buku ini berguna, daripada menguap di laptop kami masing-masing. Beberapa materi diambil dari materi yang diajarkan oleh Bapak Matradji dan Bapak Yerri Susatio di Departemen Teknik Fisika ITS. Kontributor tulisan per-babnya adalah:  
\begin{itemize}
    \item Bab 1 Persamaan Garis Lurus oleh B.T. Atmaja 
    \item Bab 2 Fungsi oleh B.T. Atmaja 
    \item Bab 3 Pendekatan Akar oleh B.T. Atmaja 
    \item Bab 4 Interpolasi oleh M.N. Farid  
    \item Bab 5 Integral Kalkulus oleh D. Prananto 
    \item Bab 6 Deret Fourier oleh D. Prananto
    \item Bab 7 Analisa Sistem oleh D. Prananto  
\end{itemize}
Bab 1-3 diambil dari materi yang diajarkan oleh Bapak Yerri Susatio, Bab 4 disalin dari tulisan tangan Bapak Matradji, Bab 5-7 diambil dari Buku Sinyal \& Sistem karya Oppenheim dan Willsky. Fail mentah (\LaTeX) dari buku ini bisa diakses di \url{https://github.com/bagustris/matrek}.  

Kami menyadari bahwa buku ini masih jauh dari sempurna. Oleh karena itu, masukan dan kritik yang membangun sangat kami harapkan demi perbaikan edisi selanjutnya. Semoga buku ini bermanfaat bagi para pembaca dan memberikan sumbangsih nyata dalam dunia pendidikan teknik di Indonesia.

\vspace{1cm}
\noindent Surabaya, 2017   

\vspace{0.5cm}
\noindent \textit{Penulis}

